% !TITLE 梦微之
% !AUTHOR 白居易
% !DYNASTY 唐
% !GENRE 七言律诗
% !ORDER 49

\begin{poem}
夜来携手梦同游，晨起盈巾泪莫收。
漳浦老身三度病，咸阳宿草八回秋。
君埋泉下泥销骨，我寄人间雪满头。
阿卫韩郎相次去，夜台茫昧得知不？
\end{poem}

\begin{appreciation}
微之是元稹的字。他和白居易同年进士及第，一生志同道合，唱和极多，被后人并称"元白"。元稹先白居易九年去世，葬于咸阳。这首诗写于元稹去世八年后，白居易已年过七旬。

夜来携手梦同游，晨起盈巾泪莫收——夜里梦见和元稹携手同游，早上起来，泪水打湿了枕巾，擦也擦不干。盈巾是满巾的意思。一个"盈"字，写出了泪水之多、悲伤之深。但白居易没有写自己如何痛哭，只是平静地说"泪莫收"——收不住了，也不想收了。

漳浦老身三度病，咸阳宿草八回秋——漳（zhāng）浦是白居易当时居住的地方，他在那里已经病过三次。咸阳是元稹葬地，宿草是隔年的草，八回秋就是八年。一边是病中的老人，一边是坟上的宿草；一边是活着的煎熬，一边是死后的沉寂。两句话，写尽了生死相隔的苍凉。

君埋泉下泥销骨，我寄人间雪满头——你埋在黄泉之下，泥土销蚀着你的骨头；我寄居在人间，满头白发如雪。销（xiāo）是销蚀、消融的意思。这两句是千古名句，对仗工整而情感深挚。泉下与人间，泥与雪，销骨与满头，生死对比，触目惊心。白居易没有说"我想你"，而是说"你在土里，我在世上"，这种冷静的陈述，比任何呼号都更让人心碎。

阿卫韩郎相次去，夜台茫昧得知不——阿卫是元稹的小儿子，韩郎是元稹的女婿。相次去是一个接一个地去世了。夜台是坟墓的别称，茫昧（máng mèi）是昏暗不明的意思。不在这里读fǒu，是疑问词。白居易问：你的儿女也都相继去世了，你在那黑暗的地下，知道吗？

这个问句是绝望的。他问的是一个不可能有答案的问题。但正是这种无望的追问，写出了思念的极致。白居易晚年丧友、丧子、丧婿，亲人朋友一个个离去，只剩下他一个人"寄人间"。那个"寄"字，道尽了一个老人的孤独——人间已不是他的家，他只是暂时寄居于此，终有一天，也要去那夜台之下，与故人相见。
\end{appreciation}
